\documentclass[12pt, twoside]{book}
\usepackage[utf8]{inputenc}
\usepackage[spanish]{babel}
\usepackage[width=16cm,height=21cm]{geometry}
\usepackage{fancyhdr}
\usepackage[square,numbers]{natbib}
\usepackage[nottoc,numbib]{tocbibind}
\pagestyle{fancy}
\setlength{\headheight}{20pt}
\linespread{1.1}
\fancyhead[E]{\thepage\hfill \textit{\nouppercase{\rightmark}}}
\fancyhead[O]{\textit{\nouppercase{\leftmark}}\hfill\thepage}
\renewcommand{\headrulewidth}{0.5pt}
\cfoot{}
\usepackage{amsmath,amssymb,mathtools}
\usepackage{url}
\usepackage{tikz}
\usetikzlibrary{shapes.geometric}
\newcommand{\bD}{\mathbb{D}}
\newcommand{\bR}{\mathbb{R}}
\newcommand{\bZ}{\mathbb{Z}}
\newcommand{\bC}{\mathbb{C}}
\newcommand{\bN}{\mathbb{N}}
\newcommand{\cH}{\mathcal{H}}
\newcommand{\cA}{\mathcal{A}}
\renewcommand{\phi}{\varphi}
\newcommand{\arctanh}{\operatorname{arctanh}}
\newcommand{\sinc}{\operatorname{sinc}}
\newcommand{\eqdef}{\coloneqq}
\newcommand{\eps}{\varepsilon}
\newcommand{\de}{\delta}
\usepackage{amsthm}
\newtheorem{thm}{Teorema}
\newtheorem{prop}[thm]{Proposición}
\newtheorem{lem}[thm]{Lema}
\newtheorem{conj}[thm]{Conjetura}
\theoremstyle{definition}
\newtheorem{defn}[thm]{Definición}
\newtheorem{cor}[thm]{Corolario}
\newtheorem{ex}[thm]{Ejemplo}



\usepackage{xcolor}
\newenvironment{egornote}{\color{blue}}{}

\title{ }
\author{ \\
[1cm]{Directores de tesis:
}}
\date{July 2021}

\begin{document}

   \begin{titlepage}
        \thispagestyle{empty}
        \begin{minipage}[c][0.17\textheight][c]{0.25\textwidth}
            \begin{center}
                \includegraphics[width=2.5cm, height=3.5cm]{imagenes/Logo_Instituto_Politécnico_Nacional.png}
            \end{center}
        \end{minipage}
        \begin{minipage}[c][0.195\textheight][t]{0.75\textwidth}
            \begin{center}
                \vspace{0.3cm}
                \textsc{\large INSTITUTO POLITÉCNICO NACIONAL}\\[0.5cm]
                \vspace{0.3cm}
                \hrule height2.5pt
                \vspace{.2cm}
                \hrule height1pt
                \vspace{.8cm}
                \textsc{ESCUELA SUPERIOR DE FÍSICA Y MATEMÁTICAS }\\[0.5cm] %
            \end{center}
        \end{minipage}

        \begin{minipage}[c][0.81\textheight][t]{0.25\textwidth}
            \vspace*{5mm}
            \begin{center}
                \hskip2.0mm
                \vrule width1pt height13cm 
                \vspace{5mm}
                \hskip2pt
                \vrule width2.5pt height13cm
                \hskip2mm
                \vrule width1pt height13cm \\
                \vspace{5mm}
                \includegraphics[height=4.0cm]{imagenes/ESCUDO_ESFM.png}
            \end{center}
        \end{minipage}
        \begin{minipage}[c][0.81\textheight][t]{0.75\textwidth}
            \begin{center}
                \vspace{1cm}

                {\large\scshape Aquí va el título}\\[.2in]

                \vspace{2cm}            

                \textsc{\LARGE T\hspace{1.5cm}E\hspace{1.5cm}S\hspace{1.5cm}I\hspace{1.5cm}S}\\[0.5cm]
                \textsc{\large que para obtener el grado de:}\\[0.5cm]
                \textsc{\large Maestro en Ciencias Fisicomatemáticas}\\[0.5cm]
                \textsc{\large presenta:}\\[0.5cm]
                \textsc{\large {Tu nombre}}\\[2cm]          

                \vspace{0.5cm}

                {\large\scshape Directores de tesis:\\[0.3cm] {Director 1 \\[0.2cm]
        
                Director 2}}\\[.2in]

                \vspace{0.5cm}

                \large{Ciudad de México,}{ }{2022}
            \end{center}
        \end{minipage}
    \end{titlepage}

\section*{Resumen}




\clearpage
\section*{Abstract}


\tableofcontents

\chapter*{Introducci\'on}
\addcontentsline{toc}{chapter}{Introducción}
\input{Capitulos/c0-introduccion}

\chapter{Capitulo 1}
\input{Capitulos/c1}

\chapter{Capitulo 2}
\input{Capitulos/c2}

\chapter{Capitulo 3}
\input{Capitulos/c3}

\chapter{Capitulo 4}
\input{Capitulos/c4}

\bibliographystyle{abbrvnat}
\bibliography{sample}
\end{document}
